\begin{textsample}{POS Dim 1 – 2023\_c4 – Score -2.00 – 2023\_c4/enade\_2023\_eng\_el\_ca14304256.txt}  \label{ex:f1_pos_191}
CAPÍTULO 3 PERCEPÇÃO DA PROVA Neste capítulo, constam as percepções dos concluintes da área de Engenharia Elétrica sobre a prova aplicada.
Essas percepções foram mensuradas por meio de nove questões que avaliaram desde o grau de dificuldade da prova até o tempo gasto para concluí -la.
As percepções sobre a prova foram relacionadas ao desempenho dos estudantes, à grande região de funcionamento do curso, à categoria administrativa e ao tipo de organização acadêmica da IES.
O questionário de percepção da prova encontra -se ao final do Anexo VII, no qual está a reprodução do exame.
Na apresentação dos dados relativos às nove questões sobre as percepções a respeito da prova, o desempenho dos estudantes foi classificado em quatro quartos.
Para tanto, esse desempenho foi ordenado de forma ascendente.
O percentil 25, P25, também conhecido como primeiro quartil, é a nota de desempenho que deixa um quarto ( 25 \% ) dos valores abaixo e três quartos acima.
O quarto inferior de desempenho é composto pelas notas abaixo do primeiro quartil.
Já o percentil 75, P75, também conhecido como terceiro quartil, é o valor para o qual há três quartos ( 75 \% ) dos dados abaixo e um quarto acima dele.
O quarto superior de desempenho é composto pelas notas iguais ou acima do terceiro quartil.
O percentil 50, P50, também conhecido como mediana, é o valor que divide as notas em dois conjuntos de igual tamanho.
O 2.º quarto inclui os valores entre o primeiro quartil ( P25 ) e a mediana.
O 3.º quarto contém os valores entre a mediana ( P50 ) e o terceiro quartil ( P75 ).
Vale ressaltar que percentis, quartis e medianas são pontos que não obrigatoriamente pertencem ao conjunto original de dados, ao passo que os quartos são subconjuntos dos dados originais.
Na Figura 3.1, apresenta -se uma ilustração dos quatro quartos descritos.
A seguir, serão apresentados gráficos com resultados selecionados, relativos às nove questões 78 Relatório Síntese de Área – ENADE/2023 avaliadas por grupos de estudantes.
Nas barras dos gráficos, apresenta -se o percentual de estudantes que assinalaram uma das opções ou a soma das porcentagens daqueles que assinalaram duas ( ou três ) delas.
Por exemplo, para as questões 1 e 2, nos gráficos é apresentada a porcentagem total de participantes que assinalaram as opções ( D ) “ Difícil ” e ( E ) “ Muito difícil ”.
Em cada barra, foram assinalados também os extremos do intervalo de confiança de 95 \% como linhas verticais unidas por uma linha horizontal na forma da letra H maiúscula.
O estimador de um parâmetro com um certo nível de confiança ( por exemplo, 95,0 \% ) deve conter o parâmetro no intervalo de confiança em 95 \% das vezes.
Na comparação entre os estimadores dos parâmetros de duas classes de uma dada categoria ( por exemplo : Norte e Nordeste nas grandes regiões, de primeiro e último quarto dentro de desempenho, de IES públicas e privadas, ou de universidades e faculdades ), associados aos seus respectivos intervalos de confiança, diz -se que não há diferença estatisticamente significativa entre os parâmetros das duas categorias quando há interseção entre os intervalos de confiança, e que há diferença, se os intervalos de confiança são disjuntos ( vide Glossário ).
Nas tabelas do Anexo V, são apresentados os valores absolutos e a distribuição percentual das alternativas válidas das nove questões, segundo o mesmo recorte de desempenho dos estudantes e da grande região de funcionamento do curso, de categoria administrativa e do tipo de organização acadêmica da IES. 3.1 GRAU DE DIFICULDADE DA PROVA Nesta seção, são mostradas as percepções dos estudantes quanto ao grau de dificuldade da prova do Enade 2023 que foram analisadas considerando sua vinculação ao componente de Formação Geral e ao componente de Conhecimento Específico. 3.1.1 COMPONENTE DE FORMAÇÃO GERAL Ao avaliarem “ Qual o grau de dificuldade das questões de Formação Geral? ” ( Questão 6 ), 30,7 \% do grupo de inscritos e presentes optaram pelas alternativas “ Difícil ” ou “ Muito difícil ”.
Entretanto, para 52,7 \% dos estudantes, o componente de Formação Geral da prova foi avaliado como tendo grau de dificuldade “ Médio ” ( Gráfico 3.1, Gráfico 3.2 e Anexo V, Tabela V.1 ).
De acordo com os dados do Gráfico 3.1, o percentual de estudantes que consideraram a prova “ Difícil ” ou “ Muito difícil ” foi maior na região Sul, onde a proporção foi de 36,0 \%, enquanto a menor proporção foi identificada na região Nordeste, com 25,1 \%.
O maior intervalo foi observado na região Centro-Oeste, com variação de 7,2 \% entre o mínimo e máximo.
Além disso, é possível observar que as diferenças entre as regiões Sul e Centro-Oeste em relação às regiões Nordeste, Norte e Sudeste são estatisticamente significativas.
Nas grandes regiões, a proporção de estudantes presentes que avaliaram o componente de Formação Geral como tendo grau de dificuldade “ Médio ” foi de 48,4 \% na região Centro-Oeste, 55,9 \% na região Nordeste, 58,6 \% na região Norte, 52,8 \% na região Sudeste e 49,6 \% na região Sul.
O Gráfico 3.2 apresenta mais resultados sobre a questão tratada no gráfico anterior, porém por quartos de desempenho.
Os dados do Gráfico 3.2 mostram que o percentual de estudantes que consideraram o componente de Formação Geral da prova “ Difícil ” ou “ Muito difícil ” foi decrescente em função dos quartos de desempenho : 40,0 \% no 1.º quarto e 21,4 \% no 4.º quarto, o grupo de melhor desempenho na prova.
Nos quartos de desempenho intermediários, a proporção de estudantes que consideraram a prova “ Difícil ” ou “ Muito difícil ” foi de 33,6 \% no 2.º quarto, e de 28,3 \% no 3.º quarto.
Há diferenças estatisticamente significativas entre os quartos de desempenho.
Para todos os quartos de desempenho, a alternativa modal para essa pergunta foi a “ Médio ”, com 51,6 \% e 52,5 \% dos respondentes nos quartos extremos, 1.º e 4.º, respectivamente.
O Gráfico 3.3 apresenta outros resultados em relação à mesma questão retratada no gráfico anterior, porém por organização acadêmica das IES.
Como mostram os dados do Gráfico 3.3, o percentual de estudantes que consideraram o componente de Formação Geral da prova “ Difícil ” ou “ Muito difícil ” foi maior entre os estudantes de faculdades, alcançando 42,5 \%, enquanto entre os estudantes das CEFET/IF, organizações acadêmicas com as menores incidências, os percentuais alcançaram 26,4 \%.
É possível observar que há diferenças estatisticamente significativas entre as organizações acadêmicas dos dois extremos em relação as organizações acadêmicas das posições centrais do gráfico.
Nas organizações acadêmicas, a proporção de estudantes presentes que avaliaram o componente de Formação Geral como de grau de dificuldade “ Médio ” foi de 53,8 \% nas universidades, 49,7 \% nos centros universitários, 48,9 \% nas faculdades e 55,8 \% nos CEFET/IF ( ver também a Tabela V.2 no Anexo V ).
O Gráfico 3.4 apresenta os dados por categoria administrativa das IES.
Os dados do Gráfico 3.4 mostram que o percentual de estudantes que consideraram o componente de Formação Geral da prova “ Difícil ” ou “ Muito difícil ” foi maior entre os estudantes de cursos de IES privada ( 39,9 \% ), com diferença estatisticamente significativa da proporção registrada em IES \textbf{pública} ( 19,7 \% ).
Já a proporção de estudantes presentes que avaliaram o componente de Formação Geral como de grau de dificuldade “ Médio ” foi de 57,0 \% entre os participantes de IES públicas e de 49,1 \% entre os de IES privadas ( ver também a Tabela V.2 no Anexo V ). 3.1.2 COMPONENTE DE CONHECIMENTO ESPECÍFICO Ao responderem à questão “ Qual o grau de dificuldade das questões do componente Específico? ”, 61,2 \% do grupo de estudantes classificaram -na como “ Difícil ” ou “ Muito difícil ”.
Além disso, o componente modal de componente de Conhecimento Específico da prova foi considerado com grau de dificuldade “ Difícil ” por 51,4 \% dos estudantes ( Gráfico 3.5, Gráfico 3.6, e, no Anexo V, a Tabela V.3 ).
O Gráfico 3.5 apresenta as respostas obtidas para a Questão 8 : “ Qual o grau de dificuldade das questões do componente Específico? ”.
De acordo com os dados do Gráfico 3.5, a análise das respostas dos estudantes quanto ao grau de dificuldade do componente de Conhecimento Específico da prova, agregado por grande região, indica que a região Sul apresentou o maior percentual de estudantes que consideraram o componente de Conhecimento Específico da prova “ Difícil ” ou “ Muito difícil ”, com 64,3 \%, seguida da região Centro-Oeste, com 63,7 \%.
As diferenças entre a região Nordeste e as regiões Centro-Oeste, Sudeste e Sul são estatisticamente significativas.
Os estudantes que classificaram o grau de dificuldade como “ Médio ”, no componente de Conhecimento Específico, atingiram um percentual que variou de 32,8 \% a 41,4 \%, para as regiões Centro-Oeste e Nordeste, respectivamente.
O Gráfico 3.6 apresenta mais resultados sobre a questão tratada no gráfico anterior, porém por quartos de desempenho.
Os dados do Gráfico 3.6 mostram que o percentual de estudantes que consideraram o componente de Conhecimento Específico da prova “ Difícil ” ou “ Muito difícil ” foi decrescente em função dos quartos de desempenho, com : 64,7 \% no 1.º quarto e 54,8 \% no 4.º quarto, o grupo de melhor desempenho na prova.
Nos quartos de desempenho intermediários, a proporção de estudantes que consideraram a prova “ Difícil ” ou “ Muito difícil ” foi de 64,6 \% no 2.º quarto e de 62,7 \% no 3.º quarto.
Observa -se que há diferenças estatisticamente significativas dos resultados entre o último quarto e os demais quartos de desempenho.
Já os que responderam que o grau de dificuldade das questões do componente de Conhecimento Específico da prova foi “ Médio ” atingiram percentagens que variaram de 32,1 \%, no 1.º quarto, a 41,6 \%, no 4.º quarto.
O Gráfico 3.7 apresenta outros resultados em relação à mesma questão retratada no gráfico anterior, porém por organização acadêmica das IES.
Como mostram os dados do Gráfico 3.7, o percentual de estudantes que consideraram o componente de Conhecimento Específico da prova “ Difícil ” ou “ Muito difícil ” foi maior entre os estudantes de faculdades ( 69,2 \% ), enquanto entre os estudantes das universidades, a organização acadêmica com a menor incidência, o percentual alcançou 57,1 \%.
Além disso, há diferenças estatisticamente significativas entre as universidades e as demais organizações acadêmicas.
Já as proporções de estudantes que consideraram o componente de Conhecimento Específico como de grau de dificuldade “ Médio ” estiveram entre 27,3 \% nas faculdades, e 39,2 \% nas universidades ( ver também Tabela V.4, no Anexo V ).
O Gráfico 3.8 apresenta os dados por categoria administrativa das IES.
Os dados do Gráfico 3.8 mostram que o percentual de estudantes que consideraram o componente de Conhecimento Específico da prova “ Difícil ” ou “ Muito difícil ” foi maior entre os estudantes de cursos de IES privada ( 67,7 \% ), uma diferença estatisticamente significativa da proporção em relação às IES pública ( 53,4 \% ).
Já a proporção de presentes na prova que consideraram este componente de grau de dificuldade “ Médio ” foi de 43,0 \% para os estudantes de IES públicas e 29,5 \% para os de IES privadas ( ver Tabela V.4, Anexo V ). 3.2 EXTENSÃO DA PROVA EM RELAÇÃO AO TEMPO TOTAL Indagados quanto à extensão da prova, em relação ao tempo total oferecido para a sua resolução ( Questão 2 ), os estudantes apontaram, com maior incidência, a alternativa que considerava a extensão “ Adequada ”, em todas as agregações consideradas ( Gráfico 3.9, Gráfico 3.10, e, no Anexo V, a Tabela V.5 ).
O Gráfico 3.9 apresenta as respostas obtidas para a Questão 2 : “ Em relação ao tempo total de aplicação, você considera que a prova foi longa ou muito longa? ” Em relação aos dados do Gráfico 3.9, entre as grandes regiões, os que avaliaram a prova como “ longa ” ou “ muito longa ”, em relação ao tempo total destinado a sua resolução, resultaram em porcentagens que variam de 26,1 \%, na região Norte, até 33,8 \%, na região Sudeste.
É possível observar que há diferenças estatisticamente significativas entre as regiões Sul e Sudeste em relação as regiões Norte e Nordeste.
O percentual de estudantes que responderam ser a extensão da prova “ Adequada ” foi de 60,9 \%.
Já 31,7 \% dos inscritos presentes consideraram que a prova foi “ longa ” ou “ muito longa ”, e 7,3 \% a avaliaram como “ Curta ” ou “ Muito curta ”.
O Gráfico 3.10 traz mais resultados sobre a questão tratada no gráfico anterior, porém por quartos de desempenho.
No Gráfico 3.10, observa -se que a proporção de estudantes que consideraram a prova “ longa ” ou “ muito longa ” foi 33,0 \% no 4.º quarto — grupo de estudantes com melhor desempenho — e também 33,0 \% no 1.º quarto.
Pode -se constatar que não há diferença estatisticamente significativa entre o 1.º quarto e os demais quartos de desempenho.
Considerando -se o desempenho dos estudantes, nota -se que 60,9 \% consideraram a extensão da prova “ Adequada ”, no quarto de desempenho inferior, e 57,7 \%, no de melhor desempenho ( quarto superior ).
Nos quartos intermediários, essa proporção foi de 63,5 \% no 2.º quarto e de 61,6 \% no 3.º.
O Gráfico 3.11 traz outros resultados também em relação à questão retratada no gráfico anterior, porém por organização acadêmica das IES.
De acordo com os dados do Gráfico 3.11, considerando -se a extensão da prova, o percentual de estudantes que avaliaram a prova como “ longa ” ou “ muito longa ”, em relação ao tempo total destinado à sua resolução, foi maior entre os estudantes de universidades ( 33,2 \% ), enquanto os de centros universitários, a organização acadêmica com a menor incidência, esse índice atingiu 28,1 \%.
Pode -se constatar que há diferença estatisticamente significativa entre os centros universitários em relação às universidades e os CEFET/IF.
A proporção de presentes, na prova do Enade 2023, que consideraram a sua extensão “ Adequada ”, em relação ao tempo total destinado à sua resolução, esteve entre 57,9 \% nos CEFET/IF e 64,7 \% nos centros universitários ( ver Anexo V, Tabela V.6 ).
O Gráfico 3.12 apresenta os dados por categoria administrativa das IES.
Os dados do Gráfico 3.12 mostram que o percentual de estudantes que consideraram a prova “ longa ” ou “ muito longa ”, em relação ao tempo total destinado à sua resolução, foi maior entre os estudantes de cursos de IES \textbf{pública} ( 33,1 \% ), sem diferença estatisticamente significativa da proporção registrada em IES privada ( 30,6 \% ).
Já a proporção de estudantes presentes na prova que consideraram a extensão “ Adequada ” foi de 59,1 \%, entre os estudantes de IES públicas, e de 62,4 \%, entre os de IES privadas ( ver a Tabela V.6 no Anexo V ). 3.3 COMPREENSÃO DOS ENUNCIADOS DAS QUESTÕES Também organizados em relação ao componente de Formação Geral e ao componente de Conhecimento Específico, os dados apresentados nesta seção dizem respeito à forma como os estudantes compreenderam os enunciados das questões da prova, considerando as formulações expressas nesses mesmos enunciados. 3.3.1 COMPONENTE DE FORMAÇÃO GERAL Em relação aos enunciados das questões do componente de Formação Geral, as opiniões foram positivas, uma vez que 78,2 \% dos estudantes avaliados consideraram todos ou a maioria dos enunciados das questões “ Compreensíveis e objetivos ” ( Gráfico 3.13, Gráfico 3.14 e, no Anexo V, a Tabela V.7 ).
O Gráfico 3.13 apresenta as respostas obtidas para a Questão 7 : “ Os enunciados das questões de Formação Geral estavam compreensíveis e objetivos? ”.
Os resultados apresentados no Gráfico 3.13 mostram que, na análise regional, a porcentagem de estudantes que avaliaram que todos ou a maioria dos enunciados das questões do componente de Formação Geral estavam “ compreensíveis e objetivos ” variou em 76,3 \% ( menor porcentagem ), na região Sul e 80,4 \% ( maior porcentagem ), na região Nordeste, não sendo observadas diferenças estatisticamente significativas entre as proporções das regiões.
O Gráfico 3.14 apresenta mais resultados sobre a questão tratada no gráfico anterior, porém por quartos de desempenho.
De acordo com os dados do Gráfico 3.14, segundo o desempenho, observa -se que a proporção dos estudantes que emitiram essa opinião cresce, conforme o desempenho aumenta, com diferenças estatisticamente significativas entre todos os quartos de desempenho.
No quarto superior, a clareza e a objetividade de todos ou da maioria dos enunciados das questões foram percebidas por 83,9 \%.
Já no 1.º quarto de desempenho, todos os enunciados das questões estavam “ compreensíveis e objetivos ” para 69,9 \% dos estudantes.
O Gráfico 3.15 traz outros resultados em relação à mesma questão retratada no gráfico anterior, porém por organização acadêmica das IES.
Os dados do Gráfico 3.15 mostram que a proporção dos presentes que consideraram que todos ou a maioria dos enunciados das questões do componente de Formação Geral estavam “ compreensíveis e objetivos ” foi maior entre os estudantes de CEFET/IF, chegando a 79,6 \%.
Entre os estudantes de universidades, o índice atingiu 78,8 \%, dos centros universitários, 77,0 \%, e de faculdades, 75,5 \%.
É possível observar que as diferenças entre as organizações acadêmicas não são estatisticamente significativas ( ver também Tabela V.8 no Anexo V ).
O Gráfico 3.16 apresenta os dados por categoria administrativa das IES.
O Gráfico 3.16 mostra que o percentual de estudantes que consideraram que todos ou a maioria dos enunciados das questões do componente de Formação Geral estavam “ compreensíveis e objetivos ” foi maior entre os estudantes de cursos de IES \textbf{pública} ( 81,3 \% ), com diferença estatisticamente significativa da proporção registrada em IES privada ( 75,7 \% ). 3.3.2 COMPONENTE DE CONHECIMENTO ESPECÍFICO Em relação aos enunciados das questões do componente de Conhecimento Específico da prova, para 75,7 \% dos estudantes avaliados da área Engenharia Elétrica, a clareza e a objetividade ( Questão 9 ) estavam presentes em todas ou na maioria das questões ( Gráfico 3.17 a Gráfico 4.20 e, no Anexo V, as Tabelas V.9 e V.10 ).
O Gráfico 3.17 apresenta as respostas obtidas para a Questão 9 : “ Os enunciados das questões do componente Específico estavam compreensíveis e objetivos? ”.
A maioria dos estudantes de todas as grandes regiões brasileiras considerou “ compreensíveis e objetivos ” todos ou a maioria dos enunciados das questões do componente de Conhecimento Específico da prova.
A maior porcentagem ocorreu na região Nordeste ( 76,8 \% ).
As diferenças entre as regiões não são estatisticamente significativas.
O Gráfico 3.18 apresenta mais resultados sobre a questão tratada no gráfico anterior, porém por quartos de desempenho.
Os dados do Gráfico 3.18 mostram que a proporção de estudantes que consideraram os enunciados das questões “ compreensíveis e objetivos ” apresenta uma tendência crescente em relação ao aumento de desempenho, com percentual mais elevado no quarto superior ( 81,9 \% ), se comparado ao quarto inferior ( 69,5 \% ).
As diferenças entre todos os quartos de desempenho são estatisticamente significativas, exceto entre o primeiro e o segundo.
O Gráfico 3.19 apresenta outros resultados em relação à mesma questão retratada no gráfico anterior, porém por organização acadêmica das IES.
Os dados do Gráfico 3.19 mostram que a proporção dos presentes que consideraram que todos ou a maioria dos enunciados das questões do componente de Conhecimento Específico estavam “ compreensíveis e objetivos ” foi maior entre os estudantes de universidades ( 76,6 \% ).
Entre os estudantes de CEFET/IF, centros universitários e faculdades, essa proporção atingiu 76,3 \%, 74,2 \% e 71,7 \%, respectivamente.
É possível observar que as diferenças entre as organizações acadêmicas não são estatisticamente significativas ( ver também Tabela V.10, no Anexo V ).
O Gráfico 3.20 apresenta os dados por categoria administrativa das IES.
Os dados do Gráfico 3.20 mostram que o percentual de estudantes que consideraram que todos ou a maioria dos enunciados das questões do componente de Conhecimento Específico estavam “ compreensíveis e objetivos ” foi maior entre os estudantes de cursos de IES \textbf{pública} ( 80,1 \% ) do que entre os de IES privada ( 71,9 \% ), sem diferença estatisticamente significativa. 3.4 SUFICIÊNCIA DAS INFORMAÇÕES/INSTRUÇÕES FORNECIDAS Em relação às informações/instruções fornecidas para a resolução das questões ( Questão 3 ), constatou -se que 79,5 \% dos respondentes da área de Engenharia Elétrica de todo o Brasil afirmaram que essas eram até excessivas ou suficientes em todas ou na maioria das questões ( Gráfico 3.21, Gráfico 3.22 e, no Anexo V, a Tabela V.11 ).
O Gráfico 3.21 apresenta as respostas obtidas para a Questão 3 : “ As informações/instruções fornecidas para a resolução das questões foram suficientes para resolvê -las? ”.
De acordo com os dados do Gráfico 3.21, quanto à distribuição de respondentes pelas grandes regiões, observa -se que a proporção de estudantes que consideraram as informações/instruções fornecidas excessivas ou suficientes em todas ou na maioria das questões foi, no mínimo, 78,1 \% ( região Sul ), alcançando até 80,8 \%, na região Nordeste.
As diferenças entre as regiões não são estatisticamente significativas.
O Gráfico 3.22 apresenta resultados sobre a questão tratada no gráfico anterior, porém por quartos de desempenho.
De acordo com os dados do Gráfico 3.22, tendo em vista o desempenho dos participantes, observa -se uma tendência crescente das proporções de participantes que avaliaram as informações/instruções como excessivas ou suficientes em todas ou na maioria das questões, sendo que todos os quartos de desempenho apresentaram diferença estatisticamente significativa.
O percentual foi mais elevado no quarto superior ( 86,0 \% ), com percentual superior ao da média nacional ( 79,5 \% ).
No quarto inferior, a suficiência das informações/instruções declarada como até excessiva, em todas ou na maioria das questões foi percebida por 73,8 \% dos respondentes.
Já nos quartos intermediários, 2.º e 3.º, as informações/instruções foram consideradas até excessivas ou suficientes em todas ou na maioria das questões por 77,1 \% e 80,9 \% dos participantes, respectivamente.
O Gráfico 3.23 apresenta outros resultados também em relação à essa questão retratada no gráfico anterior, porém por organização acadêmica das IES.
De acordo com os dados do Gráfico 3.23, considerando -se as informações/instruções fornecidas para a resolução das questões ( Questão 3 ), a proporção que considerou que essas eram excessivas ou suficientes em todas ou na maioria das questões foi maior entre os estudantes de CEFET/IF, totalizando 81,8 \%.
As faculdades foram as organizações acadêmicas em que houve a menor proporção : 75,4 \%.
É possível observar que há diferenças estatisticamente significativas entre as organizações acadêmicas dos dois extremos em relação as organizações acadêmicas das posições centrais do gráfico.
O Gráfico 3.24 apresenta os dados por categoria administrativa das IES.
De acordo com os dados do Gráfico 3.24, o percentual de estudantes que consideraram que as informações/instruções fornecidas para a resolução das questões eram até excessivas ou suficientes em todas ou na maioria das questões foi maior entre os estudantes de cursos de IES \textbf{pública} ( 83,9 \% ) em relação aos de IES privada ( 75,9 \% ), portanto, com diferença estatisticamente significativa entre as proporções registradas. 3.5 DIFICULDADE ENCONTRADA AO RESPONDER À PROVA Quando indagados sobre as dificuldades com as quais se depararam ao responder à prova ( Questão 4 ), 18,8 \% dos estudantes apontaram o “ Desconhecimento do conteúdo ”.
Para 45,8 \%, a “ Forma diferente de abordagem do conteúdo ” foi indicada como dificuldade.
Já a “ Falta de motivação para fazer a prova ” foi a dificuldade apontada por 19,1 \% dos respondentes.
Considerando -se todo o Brasil, 12,0 \% dos respondentes afirmaram que não tiveram qualquer tipo de dificuldade para responder à prova ( Gráficos de 3.25 a 3.28 e Tabelas V.13 e V.14, ambas no Anexo V ).
O Gráfico 3.25 apresenta as respostas obtidas sobre o “ Desconhecimento do conteúdo ” como a principal “ dificuldade ao responder à prova ”.
De acordo com os dados do Gráfico 3.25, os percentuais de inscritos e presentes que apontaram o “ Desconhecimento do conteúdo ” como “ dificuldade ao responder à prova ” variaram de 17,2 \%, na região Centro-Oeste, a 23,5 \%, na região Norte, sendo estatisticamente significativa a diferença entre essas duas regiões.
O tipo de dificuldade “ Forma diferente de abordagem do conteúdo ” foi a escolha modal dos estudantes, com percentuais que variaram de 42,4 \% ( região Nordeste ) a 47,7 \% ( região Norte ).
Em contrapartida, o percentual de estudantes que citaram como dificuldade “ Espaço insuficiente para responder às questões ” variou de 3,8 \% ( região Sul ) a 5,5 \% ( região Centro-Oeste ).
O Gráfico 3.26 apresenta os resultados sobre o ponto tratado no gráfico anterior, porém por quartos de desempenho.
De acordo com os dados do Gráfico 3.26, em relação aos quartos de desempenho, o “ Desconhecimento do conteúdo ” foi a opção escolhida por 16,4 \% dos estudantes do quarto inferior e por 22,7 \% do quarto superior, sendo observadas diferenças estatisticamente significativas entre o último quarto e os demais quartos de desempenho.
Levando -se em conta a alternativa modal, “ Forma diferente de abordagem do conteúdo ”, observa -se a seguinte distribuição de estudantes nos quartos de desempenho : no 1.º quarto, 50,7 \% ; no 2.º quarto, 48,0 \% ; no 3.º quarto, 46,0 \% ; e, no 4.º quarto, 38,7 \% escolheram essa alternativa.
As proporções são decrescentes com o desempenho.
O Gráfico 3.27 apresenta outros resultados em relação à mesma questão retratada no gráfico anterior, porém por organização acadêmica das IES.
De acordo com os dados do Gráfico 3.27, na análise por organização acadêmica, o percentual de inscritos e presentes que apontaram o “ Desconhecimento do conteúdo ” como “ dificuldade ao responder à prova ” não superou 19,1 \%.
Os percentuais variaram de 18,3 \%, nos CEFET/IF e nas faculdades a 19,1 \%, nas universidades.
É possível observar que as diferenças entre as organizações acadêmicas não são estatisticamente significativas.
O tipo de dificuldade “ Forma diferente de abordagem do conteúdo ” foi a escolha modal dos estudantes, com percentuais que variaram de 45,6 \% ( CEFET/IF ) a 55,7 \% ( faculdades ).
Em contrapartida, o percentual de estudantes que citaram a dificuldade “ Espaço insuficiente para responder às questões ” variou de 3,2 \% ( faculdades ) a 4,9 \% ( CEFET/IF ) ( ver Tabela V.14 no Anexo V ).
O Gráfico 3.28 apresenta os dados por categoria administrativa das IES.
De acordo com os dados do Gráfico 3.28, em relação às categorias administrativas, o “ Desconhecimento do conteúdo ” foi a opção escolhida por 20,1 \% dos estudantes nas IES públicas e 17,6 \% nas IES privadas, portanto, a diferença entre as duas categorias é estatisticamente significativa.
As alternativas modais entre os estudantes, quando agregados pelas categorias administrativas, para a dificuldade encontrada, foram as categorias “ Forma diferente de abordagem do conteúdo e Forma diferente de abordagem do conteúdo ”, com 39,2 \%, entre os estudantes de IES públicas, e 51,3 \%, entre os de IES privadas, respectivamente ( ver Tabela V.14 no Anexo V ). 3.6 CONTEÚDOS DAS QUESTÕES OBJETIVAS DA PROVA Quando analisados os conteúdos das questões objetivas da prova ( Questão 5 ), verifica -se que um percentual pequeno dos estudantes avaliados, 3,8 \%, afirmou que “ Não estudou ainda a maioria desses conteúdos ” ( Gráficos 3.29 a 4.32 e, no Anexo V, Tabelas V.15 e V.16 ).
A maior parte ( 51,0 \% ) afirmou que “ Estudou e aprendeu muitos desses conteúdos ”.
O Gráfico 3.29 apresenta as respostas obtidas para a Questão 5 : “ Considerando apenas as questões objetivas da prova, você percebeu que não estudou ainda a maioria desses conteúdos? ”.
De acordo com o Gráfico 3.29, na análise por grande região, a proporção de respondentes que escolheu a opção “ Não estudou ainda a maioria desses conteúdos ” foi pequena.
Observa -se que, nas regiões Norte ( 5,0 \% ) e Sul ( 4,0 \% ), as proporções foram maiores.
Nas regiões Centro-Oeste e Nordeste, a situação foi inversa, com proporções mais baixas, correspondendo a, respectivamente, 2,8 \% e 3,6 \%.
A região Sudeste ficou com a proporção de 3,7 \%.
As diferenças observadas entre as regiões não são estatisticamente significativas.
Ainda na análise por grande região, a maior parte dos presentes afirmou que “ Estudou e aprendeu muitos desses conteúdos ”, com proporções que variam entre 45,8 \%, na região Norte, e 52,1 \%, na região Nordeste.
O Gráfico 3.30 apresenta os resultados sobre a questão tratada no gráfico anterior, porém por quartos de desempenho.
De acordo com os dados do Gráfico 3.30, considerando -se separadamente as opiniões dos estudantes dos quatro quartos de desempenho, observa -se que, no quarto inferior, 7,6 \% ofereceram como resposta que “ Não estudou ainda a maioria desses conteúdos ”, sendo 1,2 \% os do quarto superior com a mesma resposta.
As diferenças entre os estudantes que optaram por esse motivo de dificuldade são estatisticamente significativas entre todos os quartos de desempenho.
Levando -se em conta o quarto superior, a maioria dos estudantes, 59,7 \%, afirmou que “ Estudou e aprendeu muitos desses conteúdos ”.
No outro extremo, no 1.º quarto, 43,6 \% optaram pela mesma categoria.
As proporções são crescentes com o desempenho.
O Gráfico 3.31 apresenta outros resultados também acerca dessa questão retratada no gráfico anterior, porém por organização acadêmica das IES.
De acordo com os dados do Gráfico 3.31, na análise por organização acadêmica, a proporção de respondentes que escolheu a opção “ Não estudou ainda a maioria desses conteúdos ” foi pequena.
Observa -se que as proporções nas universidades ( 3,2 \% ) e nos CEFET/IF ( 2,2 \% ) foram menores.
Nas faculdades a proporção foi mais alta, com 6,4 \% dos respondentes.
É possível observar que há diferenças estatisticamente significativas entre as organizações acadêmicas dos dois extremos em relação às organizações acadêmicas das posições centrais do gráfico.
A maior parte dos estudantes presentes afirmou que “ Estudou e aprendeu muitos desses conteúdos ”, com proporções que variam de 45,3 \% nas faculdades a 53,6 \% nos CEFET/IF ( ver também Tabela V.16, no Anexo V ).
O Gráfico 3.32 apresenta os dados por categoria administrativa das IES.
De acordo com os dados do Gráfico 3.32, em relação às categorias administrativas, observa -se que, nas IES públicas, 2,0 \% ofereceram como resposta que “ Não estudou ainda a maioria desses conteúdos ”, sendo 5,4 \% nas IES privadas com a mesma resposta.
A diferença entre os estudantes que optaram por esse nível de aprendizado nas duas categorias é estatisticamente significativa.
Considerando -se as IES públicas, 55,4 \% dos estudantes marcaram a opção “ Estudou e aprendeu muitos desses conteúdos ”, sendo de 47,3 \% a proporção dos que escolheram essa resposta nas IES privadas ( ver Tabela V.16 no Anexo V ). 3.7 TEMPO GASTO PARA CONCLUIR A PROVA Ao responderem sobre o tempo de conclusão da prova ( Questão 1 ), 80,6 \% dos estudantes afirmaram ter gasto entre duas e quatro horas ( Gráfico 3.33 a Gráfico 3.36 e, no Anexo V, as Tabelas V.17 e V.18 ).
O Gráfico 3.33 apresenta as respostas obtidas para a Questão 1 : “ Qual foi o tempo gasto por você para concluir a prova? ”.
Considerando -se as cinco grandes regiões brasileiras, os maiores percentuais de estudantes que utilizaram entre duas e quatro horas para finalizar a prova foram, nas regiões Norte, Sudeste e Sul, respectivamente, de 83,7 \%, 81,1 \% e 81,1 \%.
Na região Nordeste, o percentual de estudantes que dispensaram entre duas e quatro horas para concluir a prova foi de 78,9 \% e, na região Centro-Oeste, de 75,9 \%.
As diferenças entre as regiões Nordeste e Norte são estatisticamente significativas.
O Gráfico 3.34 apresenta os resultados sobre a questão tratada no gráfico anterior, porém por quartos de desempenho.
De acordo com os dados do Gráfico 3.34, uma vez analisadas as alternativas selecionadas pelos estudantes que se situam nos diferentes quartos de desempenho, observa -se uma proporção de 80,4 \% participantes no quarto superior que declarou ter gasto entre duas e quatro horas para concluir a prova e 79,0 \% no quarto inferior.
As diferenças não são estatisticamente significativas entre todos os quartos de desempenho.
O Gráfico 3.35 apresenta outros resultados em relação à mesma questão retratada no gráfico anterior, porém por organização acadêmica das IES.
De acordo com os dados do Gráfico 3.35, na análise por organização acadêmica, os percentuais dos estudantes que utilizaram entre duas e quatro horas para finalizar a prova foram de 79,7 \% nas universidades, 81,3 \% nos centros universitários, 81,6 \% nas faculdades e 82,7 \% nos CEFET/IF.
As diferenças não são estatisticamente significativas entre as organizações acadêmicas.
O Gráfico 3.36 apresenta os dados por categoria administrativa das IES.
De acordo com os dados do Gráfico 3.36, em relação às categorias administrativas, observa -se que, nas IES públicas, 80,8 \% dos estudantes declararam ter gasto entre duas e quatro horas para concluir a prova, sendo 80,4 \% dos estudantes das IES privadas com a mesma resposta.
A diferença não é estatisticamente significativa entre os estudantes das duas categorias ( Tabela V.18, Anexo V ).
Os resultados do Enade 2023, apresentados neste capítulo, contribuem para o entendimento de que a consideração da percepção dos estudantes concluintes sobre a prova do Enade é importante tanto para a compreensão do seu desempenho como para a reflexão sistemática e constante a respeito de determinadas realidades relacionadas a esse mesmo desempenho no país.
Por outro lado, como todos os elementos do trabalho pedagógico institucional mostram -se diretamente relacionados com o desempenho dos estudantes, também ganha relevância conhecer e analisar os resultados das percepções dos coordenadores de curso sobre a mesma prova.
As informações expressas por estudantes e coordenadores sobre questões diversas associadas aos dois componentes da prova – Formação Geral e componente de Conhecimento Específico – podem agregar valor em ações institucionais voltadas para o aperfeiçoamento do instrumento avaliativo, considerando os vários elementos envolvidos em sua elaboração.
Nesse contexto, é relevante que os resultados das percepções de estudantes e coordenadores sobre a prova do Enade ganhem relevância como um elemento a mais nas análises da oferta dos cursos nas modalidades presencial e a distância, em nível nacional.
Espera -se, assim, que, em uma visão formativa de avaliação, o conhecimento e a análise desses resultados possam contribuir para a construção de alternativas institucionais que induzam à excelência da educação superior.

% matched lemmas: pública
\end{textsample}
